\documentclass[tikz,border=6mm]{standalone}
\usepackage{lua-tikz3dtools}

\begin{document}
\foreach \frame in {5,...,33} {
    \pgfmathsetmacro{\PARAM}{-1+2*\frame/39}
    \begin{tikzpicture}[scale=2]
        \useasboundingbox[scale=0.5] (-3.15,-2.35) rectangle (3.15,2.45);
        
        % This is a table of Lua objects that we can use in our namespaced keys.
        \ltdtsetobject[
            name={return "objs"},
            object={
                local view = Matrix.zyzrotation(
                    Vector:new{1.10*pi,.36*pi,pi/2}
                )
                :multiply(
                    Matrix.perspective(
                        Vector:new{0,0,-1/18,1}
                    )
                )
                local inv=view:inverse()
                local function f(x,y)
                    return -x^3-y^2+.5
                end
                local function fy(x,y)
                    return -2*y
                end
                % This is a predicate which tells me whether a point is inside a unit cube.
                local function inside(p)
                    p=p:multiply(inv)
                    return abs(p[1])<=1.000001 and
                        abs(p[2])<=1.000001 and
                        abs(p[3])<=1.000001
                end
                % Average of triangle
                local function triangle(A,B,C)
                    return inside(
                        A:hadd(B):hadd(C):hscale(1/3)
                    )
                end
                % Average of line segment
                local function segment(A,B)
                    return inside(A:hadd(B):hscale(1/2))
                end
                % This helps iterate axis dimensions
                local function edge(axis,u,a,b)
                    if axis==1 then
                        return Vector:new{u,a,b,1}
                    elseif axis==2 then
                        return Vector:new{a,u,b,1}
                    end
                    return Vector:new{a,b,u,1}
                end
                return {
                    view=view,
                    graph={
                        f=f,
                        fy=fy
                    },
                    cube={
                        min=-1,
                        max=1,
                        triangle=triangle,
                        segment=segment,
                        edge=edge,
                        curve_filter=[[
                            return objs.cube.segment(A,B)
                        ]]
                    },
                    % I thought this part was really cool.
                    % We can save our styles as Lua Strings to be able to program with TikZ keys more nicely.
                    % I don't show an example here, but this means we could basically compute unique keys per simplex.
                    styles={
                        plane=[[
                            fill=cyan!7!white,
                            fill opacity=.40,
                            draw=none
                        ]],
                        grid=[[
                            draw=blue!32!gray,
                            line width=.11pt,
                            line cap=round
                        ]],
                        tangent=[[
                            draw=black!90,
                            line width=.52pt,
                            line cap=round
                        ]],
                        graph=[[
                            fill=red!70!orange!88!white,
                            draw=red!48!black,
                            line width=.065pt,
                            line join=round
                        ]],
                        section=[[
                            draw=green!42!black,
                            line width=.46pt,
                            line cap=round
                        ]],
                        point=[[
                            fill=yellow!84!orange,
                            draw=orange!60!black,
                            line width=.05pt
                        ]],
                        cube=[[
                            draw=black!28,
                            line width=.20pt,
                            line cap=round
                        ]]
                    }
                }
            }
        ]
        \ltdtappendsolid[
            uparams={return Vector:new{objs.cube.min,objs.cube.max,2}},
            vparams={return Vector:new{objs.cube.min,objs.cube.max,2}},
            wparams={return Vector:new{objs.cube.min,objs.cube.max,2}},
            v={return Vector:new{u,v,w,1}},
            transformation={return objs.view},
            filter={return false}
        ]
        \ltdtappendsurface[
            uparams={return Vector:new{-1.4,1.4,2}},
            vparams={return Vector:new{-1.4,1.4,2}},
            v={return Vector:new{\PARAM,v,u,1}},
            transformation={return objs.view},
            fill options={return objs.styles.plane},
            filter={return objs.cube.triangle(A,B,C)},
            % This is a lot of hard to read code so please read my next comment first,
            % as it explains how it works. (see down below)
            curve={
                local curves={}
                for _,q in ipairs({-.5,0,.5}) do
                    table.insert(curves,{
                        start=Vector:new{-1.4,q,1},
                        stop=Vector:new{1.4,q,1},
                        drawoptions=objs.styles.grid,
                        filter=objs.cube.curve_filter
                    })
                    table.insert(curves,{
                        start=Vector:new{q,-1.4,1},
                        stop=Vector:new{q,1.4,1},
                        drawoptions=objs.styles.grid,
                        filter=objs.cube.curve_filter
                    })
                end
                local y=\PARAM
                local z=objs.graph.f(\PARAM,y)
                local slope=objs.graph.fy(\PARAM,y)
                % See this comment. The way this key works is you return a lua table
                % of lua tables with certain fields filled. There are exactly 4 fields, as below.
                table.insert(curves,{
                    start=Vector:new{
                        z+slope*(-1.5-y),-1.5,1
                    },
                    stop=Vector:new{
                        z+slope*(1.5-y),1.5,1
                    },
                    drawoptions=objs.styles.tangent,
                    filter=objs.cube.curve_filter
                })
                return curves
            }
        ]
        \ltdtappendsurface[
            uparams={return Vector:new{-1.35,1.35,7}},
            vparams={return Vector:new{-1.35,1.35,7}},
            adaptive={
                return {
                    tolerance=.005,
                    feature_tolerance=.002,
                    max_depth=9,
                    max_triangles=700
                }
            },
            v={return Vector:new{u,v,objs.graph.f(u,v),1}},
            transformation={return objs.view},
            fill options={return objs.styles.graph},
            filter={return objs.cube.triangle(A,B,C)},
            curve={
                return {{
                    start=Vector:new{\PARAM,-1.35,1},
                    stop=Vector:new{\PARAM,1.35,1},
                    drawoptions=objs.styles.section,
                    filter=objs.cube.curve_filter
                }}
            }
        ]
        \ltdtappendsurface[
            uparams={return Vector:new{0,tau,10}},
            vparams={return Vector:new{0,pi,6}},
            v={
                local a=\PARAM
                local p=Vector:new{
                    a,a,objs.graph.f(a,a),1
                }
                return p:hadd(
                    Vector.sphere(Vector:new{u,v,.05})
                )
            },
            transformation={return objs.view},
            fill options={return objs.styles.point}
        ]
        \foreach \a in {-1,1} {
            \foreach \b in {-1,1} {
                \foreach \axis in {1,2,3} {
                    \ltdtappendcurve[
                        uparams={return Vector:new{-1,1,2}},
                        v={return objs.cube.edge(\axis,u,\a,\b)},
                        transformation={return objs.view},
                        draw options={return objs.styles.cube}
                    ]
                }
            }
        }
        \ltdtdisplaysimplices
    \end{tikzpicture}
}
\end{document}
